\subsection{章末確認問題}

\begin{enumerate}
  \item ソフトウェア構築は、単なるコーディングとどう違うか。
  \item 複雑さを最小化することが、なぜテストや保守に影響するのか。
  \item 「変化を予期する」と「過剰に柔軟に作る」は何が違うか。
  \item 検証しやすく構築するために、コード構造で注意すべきことを三つ挙げよ。
  \item 再利用のための構築と、再利用による構築の違いを説明せよ。
  \item 構築計画で、部品の統合順序を決めることが重要な理由を説明せよ。
  \item 依存関係管理で、ライセンスと脆弱性を確認すべき理由を説明せよ。
  \item 単体テストと統合テストの違いを説明せよ。
  \item API 設計で「誤用しにくい」ことが重要な理由を説明せよ。
  \item 表明、契約による設計、防御的プログラミングの違いを説明せよ。
  \item 分散・クラウド基盤ソフトウェアで、単一マシンのソフトウェアより難しくなる点を三つ挙げよ。
  \item 性能チューニングで、測定なしに最適化してはならない理由を説明せよ。
\end{enumerate}

\begin{plainbox}
\textbf{解答の方向性}

1 は、構築がコーディング、検証、単体テスト、統合テスト、デバッグ、再利用、統合、依存関係管理などを含む点を述べる。2 は、複雑なコードほど実行経路が多く、理解とテストが難しくなる点を述べる。3 は、予期された変化への備えと、根拠のない過剰設計の違いを述べる。4 以降は、本文の「単純で読みやすいコード」「早いフィードバック」「再利用のリスク」「段階的統合」「契約」「測定に基づく改善」といった語を使って説明できる。
\end{plainbox}
